\section*{Exercise 1}
\textit{Task: Set up the repository in GitLab.} 
\newline\newline
For this task, we initialized and structured our group repository according to the project requirements. The goal was to organize all necessary components for development, documentation, and delivery in a clean and standardized way.
\begin{itemize}
    \item \textbf{src/}: Contains the source code for the project.
    \item \textbf{bin/}: Used as the output directory where the compiled .class files and .jar files are placed.
    \item \textbf{maps/}: Stores map files required by the application.
    \item \textbf{reports/}: Contains our reports, including both the LaTeX source and the compiled PDF for each assignment.
    \item \textbf{presentations/}: Contains the PDF slides of our presentations.
    \item \textbf{gradle/wrapper/}: Contains Gradle wrapper files to ensure consistent builds across systems.
\end{itemize}

\noindent In addition to these directories, the following configuration and documentation files were added:
\begin{itemize}
    \item \textbf{.gitignore}: Ensures that build artifacts and other unnecessary files are not committed to the repository.
    \item \textbf{.gitlab-ci.yml}: Defines our GitLab CI pipeline for automatic testing and building.
    \item \textbf{README.md}: Provides instructions for compiling and running the project.
    \item \textbf{build.gradle}, \textbf{settings.gradle}, \textbf{gradlew}, \textbf{gradlew.bat}: Gradle build tool files for compiling the Java code, managing dependencies, and creating .jar files.
\end{itemize}

\noindent
This setup ensures a clean workflow for development, testing, and submission.
